\section{Exact duality and scalar optimality}
\label{sec:duality}

We use the Frobenius pairing
\(\langle X,Y\rangle=\operatorname{tr}(X^\top Y)\).  The spectral and
nuclear norms are dual.  For a compact SVD \(A=U_r\Sigma_rV_r^\top\), let
\(U_0,V_0\) be orthonormal null-space bases; a zero-dimensional factor is
allowed.  The standard formula \eqref{eq:full-subdiff} will be used without a
choice of bases inside repeated positive singular spaces.

\begin{theorem}[Strong duality and the complete solution face]
\label{thm:T1}
For arbitrary \(G\in\mathbb R^{m\times n}\) and nonzero
\(\Theta\in\mathbb R^{m\times n}\),
\begin{equation}
 p^\star
 =\max_{\substack{\|\Phi\|_2\le1\\
                   \langle\Theta,\Phi\rangle=0}}
   \langle G,\Phi\rangle
 =\min_{\lambda\in\mathbb R}\|G+\lambda\Theta\|_*.
 \label{eq:strong-duality}
\end{equation}
The dual minimizer set is nonempty and compact.  If \(\lambda^\star\) is any
dual minimizer and \(A^\star=G+\lambda^\star\Theta\), then all primal
solutions are
\begin{equation}
 \mathcal P^\star
 =\partial\|A^\star\|_*\cap\Theta^\perp.
 \label{eq:solution-face}
\end{equation}
Every dual minimizer satisfies
\begin{equation}
 |\lambda^\star|
 \le \frac{2\|G\|_*}{\|\Theta\|_*}.
 \label{eq:lambda-bound}
\end{equation}
\end{theorem}

\begin{proof}
The support function of the spectral unit ball is the nuclear norm.  Since
\(\Phi=0\) is feasible and lies in the interior of the ball,
finite-dimensional Slater--Fenchel duality gives
\eqref{eq:strong-duality}.  The reverse triangle inequality
\(\|G+\lambda\Theta\|_*\ge
|\lambda|\|\Theta\|_*-\|G\|_*\) gives coercivity and, by comparison with
\(\lambda=0\), \eqref{eq:lambda-bound}.  Finally, equality in the support
function is equivalent to
\(\Phi\in\partial\|A^\star\|_*\); imposing tangency yields exactly
\eqref{eq:solution-face}.
\end{proof}

At a fixed \(\lambda\), set
\begin{equation}
 A=G+\lambda\Theta=U_r\Sigma_rV_r^\top,
 \qquad
 a=\langle\Theta,U_rV_r^\top\rangle,
 \qquad
 C=U_0^\top\Theta V_0,
 \qquad
 \beta=\|C\|_*.
 \label{eq:a-C-beta-main}
\end{equation}

\begin{theorem}[Exact scalar interval and compact-root criterion]
\label{thm:interval}
For \(\phi(\lambda)=\|G+\lambda\Theta\|_*\),
\begin{equation}
 \partial\phi(\lambda)=[a-\beta,a+\beta].
 \label{eq:scalar-interval}
\end{equation}
Consequently, at this specified \(\lambda\):
\begin{enumerate}
 \item \(\lambda\) is dual optimal if and only if \(|a|\le\beta\);
 \item the compact equation \eqref{eq:compact-root} holds if and only if
 \(a=0\);
 \item dual optimality holds but compact stationarity fails if and only if
 \(0<|a|\le\beta\).
\end{enumerate}
Globally, \eqref{eq:compact-root} has no solution if and only if
\(a(\lambda)\ne0\) for every \(\lambda\in\argmin\phi\).  If this happens,
the dual minimizer is unique and its affine matrix is rank deficient.
\end{theorem}

\begin{proof}
The affine chain rule and \eqref{eq:full-subdiff} give
\(\partial\phi(\lambda)=\{a+\langle C,K\rangle:\|K\|_2\le1\}\), whose
endpoints are \(a\pm\|C\|_*\).  This proves
\eqref{eq:scalar-interval} and the three pointwise claims.  If two minimizers
existed, convexity would make \(\phi\) constant between them; at every
interior point \(\partial\phi=\{0\}\), hence \(a=\beta=0\) and a compact root
would exist.  Thus no-root failure forces uniqueness.  Full rectangular rank
would give \(\beta=0\), and optimality would again force \(a=0\), so the
unique no-root minimizer is rank deficient.
\end{proof}

\begin{remark}[Relative-interior nondegeneracy]
\label{rem:strict-complementarity-main}
At an optimum with \(\beta>0\), \(|a|<\beta\) is exactly
\(0\in\operatorname{ri}\partial\phi(\lambda^\star)\); equivalently, the
tangency hyperplane meets the relative interior of the kernel-block spectral
unit ball.  At \(|a|=\beta\) it meets only the exposed endpoint face
\eqref{eq:endpoint-face}.  The case \(\beta=0\) must be separated:
optimality then forces \(a=0\), so \(\partial\phi(\lambda^\star)=\{0\}\)
and \(0\in\operatorname{ri}\partial\phi(\lambda^\star)\), although the
strict inequality \(|a|<\beta\) is false.  The kernel-block slice is then
the whole spectral unit ball (or the singleton empty block when a nullity
is zero), rather than a proper exposed endpoint face.  For example,
\(G=\operatorname{diag}(1,0)\) and
\(\Theta=\bigl(\begin{smallmatrix}0&1\\1&0\end{smallmatrix}\bigr)\)
give \(\phi(\lambda)=\sqrt{1+4\lambda^2}\) and a unique rank-deficient
minimizer at zero with \(a=\beta=0\).
The nondegenerate case parallels convergence and limiting-center
phenomena for degenerate semidefinite central paths
\citep{halicka2002convergence,halicka2005limiting}, while fractional algebraic
rates are documented by Basu and Mohammad-Nezhad~\cite{basu2022central}.  Our rates follow directly
from the matrix-line normal forms in \cref{sec:rates}.
\end{remark}

The compact selection is nevertheless nondecreasing.  Indeed, for
\(\lambda_1<\lambda_2\), each compact polar factor belongs to the corresponding
nuclear-norm subdifferential, and monotonicity of the convex subdifferential
graph gives \(h_{\rm c}(\lambda_1)\le h_{\rm c}(\lambda_2)\).  The obstruction
is a jump, not a reversal of monotonicity.

\begin{remark}[Exact interval bisection]
With \(L=a-\beta\) and \(U=a+\beta\), \(U<0\) moves the bracket to the
right, \(L>0\) moves it to the left, and \(L\le0\le U\) certifies a
minimizer.  The initial bracket is \eqref{eq:lambda-bound}; on termination one
applies the completion \eqref{eq:KF-main}.  Unlike compact-polar bisection,
this test cannot jump over a minimizer, but it requires a numerical-rank
decision.  The elementary convergence proof is recorded in the supplement.
\end{remark}
